% =====================================================================
%  Conferenze, GAMM Roma (Checkpoint Festival, 3a edizione) — Intervento 2
%  Dalla tesi di laurea al videogioco storico: storia, fonti d'archivio e
%  game design
%  Fonte: trascrizione automatica "Track 2". Intervento in collegamento
%  remoto (con problemi tecnici di audio) dello sviluppatore e storico
%  medievista Fabio Belsanti, seguito da tre domande dal pubblico e da
%  una premiazione.
%  Etichette con prefisso "c1-2:". Nessun pacchetto aggiuntivo.
% =====================================================================

\chapter[Dalla tesi di laurea al videogioco storico]{Dalla tesi di laurea al
videogioco storico: fonti d'archivio, compagnie di ventura e game design}
\label{cap:c1-2-videogioco-storico}

Il secondo intervento al GAMM cambia completamente registro: dopo il panel
istituzionale del capitolo~\ref{cap:c1-1-tutela-minori}, un singolo autore
racconta \textbf{la storia di un progetto lungo quasi trent'anni}, nato da una
tesi di laurea in storia medievale e diventato un videogioco strategico.
Il relatore interviene \textbf{da remoto}, con evidenti difficoltà tecniche
(l'audio della sala gli arriva male e alla fine non si riesce a condividere il
trailer con il sonoro); la moderatrice spiega che proprio per questo l'ordine del
programma è stato cambiato.

Il relatore è \textbf{Fabio Belsanti}\footnote{Il cognome è trascritto in modi
diversi (``Bersanti'', ``Pesati''). Si adotta la forma ``Belsanti''; il nome
della società è reso come ``Job Games'', ``PM Studios'' e ``Age of the Eaties'':
con ogni probabilità si tratta di \textbf{AGE of Games}, studio di Bari, ma la
forma esatta non è verificabile dalla sola trascrizione.}, cofondatore e
amministratore di uno studio indipendente nato nel \textbf{2001} ``nel profondo
sud d'Europa'', a \textbf{Bari}. Ma, avverte subito, il progetto che presenta
\textbf{precede} la sua carriera professionale: è ``una storia delle storie'',
in cui si intrecciano la conservazione del patrimonio storico-artistico italiano,
la storia vera e propria e la sua biografia personale.

Struttura del capitolo:
\begin{enumerate}
  \item la tesi del 1998 e le fonti contabili di una compagnia di ventura
        (sezioni~\ref{sec:c1-2-tesi} e~\ref{sec:c1-2-compagnia});
  \item il ``cortocircuito'' tra libri contabili e gioco di ruolo
        (sezione~\ref{sec:c1-2-cortocircuito});
  \item il sogno interrotto e la nascita dello studio
        (sezione~\ref{sec:c1-2-studio});
  \item il game design come lente sulla storia
        (sezione~\ref{sec:c1-2-lente});
  \item il primo documento di progetto (2022) e la questione
        dell'accuratezza (sezione~\ref{sec:c1-2-documento});
  \item il gioco oggi: narrativa, indicatori, gestione e battaglie
        (sezione~\ref{sec:c1-2-gioco});
  \item le domande del pubblico (sezione~\ref{sec:c1-2-domande}) e la
        premiazione (sezione~\ref{sec:c1-2-premio}).
\end{enumerate}

% ---------------------------------------------------------------------
\section{La tesi del 1998: un archivio ad Arezzo}
\label{sec:c1-2-tesi}

Nel \textbf{1998} Belsanti era ad \textbf{Arezzo}, inviato dal suo professore
(Duccio Balestracci\footnote{Trascritto ``Adolfo Balestraci''. Duccio
Balestracci è uno storico medievista dell'Università di Siena; la correzione del
nome di battesimo è un'ipotesi ragionevole.}) a preparare una tesi di laurea
presso l'\textbf{archivio della Fraternita dei Laici di Arezzo}. Di formazione è
uno \textbf{storico medievista}, con specializzazione
\textbf{economico-militare}: ha avuto ``l'enorme fortuna e privilegio'' di
studiare all'\textbf{Università di Siena} e di passare un anno in archivio.

Il professore, conoscendo le sue passioni, gli affidò le ricchissime fonti dei
\textbf{libri contabili} della compagnia di ventura di \textbf{Micheletto
Attendolo}\footnote{Trascritto ``Michele Toltecco'', ``Michele Padendolo'',
``Micheletto Attendolo''. Micheletto Attendolo (o degli Attendoli) da Cotignola
fu un condottiero del primo Quattrocento, cugino di Muzio Attendolo Sforza.},
una compagnia di soldati mercenari del Quattrocento. Il titolo della tesi
riprende il linguaggio delle fonti: \emph{Huomini et arme de lo magnifico
signore messer Micheletto degli Attendoli de' conti di Cotignola}\footnote{La
trascrizione dà ``Ovi e tarmi di lo magnifico signore, messere Michele degli
Attenborri di Conti di Cotignoli''. La ricostruzione del titolo è
approssimativa; il senso è ``uomini d'arme del magnifico signore Micheletto
Attendolo dei conti di Cotignola''.}. Fu discussa con il massimo dei voti.

\subsection{Paolo Uccello e la Battaglia di San Romano}

La copertina della tesi riporta un'immagine della \textbf{Battaglia di San
Romano} di \textbf{Paolo Uccello}, un \textbf{trittico} i cui tre pannelli sono
oggi divisi tra \textbf{Firenze, Parigi e Londra} e che, racconta Belsanti,
``il caso vuole'' sarà riunito a dicembre per la prima volta dall'Ottocento%
\footnote{La sede della riunione è trascritta ``Palazzo Pittico'' e non è
identificabile con certezza (forse gli Uffizi o Palazzo Pitti, a Firenze).}.
La ``fortuna'' dello storico è che Micheletto Attendolo, benché non sia uno dei
condottieri più celebri (come il Gattamelata o il Colleoni), \textbf{compare nel
dipinto di Paolo Uccello}. Così, oltre ai libri contabili della compagnia,
Belsanti ebbe \textbf{un'immagine, un'icona e un'estetica di riferimento} per
quello che sarebbe diventato il gioco.

\subsection{Una delle prime digitalizzazioni d'archivio}

La tesi è oggi pubblicata gratuitamente sul portale \textbf{Academia.edu}%
\footnote{Trascritto ``accademia.eu''.}, molto dopo la laurea (1999), dopo
essere stata ``rimaneggiata un po'''; è ancora molto consultata da storici e
studenti. Il motivo del successo non è solo la parte di analisi:
\begin{itemize}
  \item fu \textbf{una delle primissime digitalizzazioni} del patrimonio
        archivistico italiano: nel 1998 Belsanti usava uno scanner ``retro
        fantascientifico'', e il professor Catoni\footnote{Giuliano Catoni,
        archivista e storico senese; la trascrizione riporta solo
        ``Catoni''.}, vecchio archivista della commissione di laurea, era
        ``ammirato e stupito'' di poter vedere e ingrandire le fonti in digitale;
  \item soprattutto, la tesi contiene la \textbf{trascrizione integrale} dei
        libri contabili (per gli anni 1424, 1425, 1430 e altri), resa
        comprensibile a chiunque e scientificamente verificabile, perché ogni
        dato rimanda alla fonte. ``Raramente uno storico professionista riesce a
        chiudersi per un anno'' a fare una trascrizione completa.
\end{itemize}

Mostra la copertina di uno dei libri, con la doppia datazione in numeri romani e
arabi (\textbf{1440}), e una delle prime pagine di un libro del
\textbf{1441}: in alto l'invocazione religiosa e il nome del tesoriere,
\textbf{Francesco di Viviano d'Arezzo}\footnote{Il nome compare come ``Francesco
Vivian Caretto'', ``Francesco di Viamonda Reco'', ``Francesco di Miliano da
Rezzo'': si ricostruisce, con riserva, come ``Francesco di Viviano d'Arezzo''.},
e la registrazione di una somma versata dalla Repubblica di Venezia\footnote{``La
signoria del campo di Necro'' nella trascrizione; dal contesto (le potenze per cui
la compagnia combatteva) è plausibile Venezia, ma il passaggio è
illeggibile.}. Una disciplina specifica dello storico è la
\textbf{paleografia}: imparare a leggere la grafia di chi scriveva i documenti.

\paragraph{Una parentesi sull'IA.} Oggi, osserva, l'intelligenza artificiale
applicata alle scansioni riesce a produrre trascrizioni notevoli molto
rapidamente; nel 1998 nulla di tutto questo esisteva. Eppure negli archivi è
digitalizzata solo una piccola parte del patrimonio: ``si sta facendo, per
fortuna, ma ci si dovrebbe investire molto di più''.

% ---------------------------------------------------------------------
\section{Che cos'era una compagnia di ventura}
\label{sec:c1-2-compagnia}

Una \textbf{compagnia di ventura} del primo Rinascimento è una compagnia di
soldati mercenari organizzata in modo \textbf{verticistico}. Al centro c'è il
\textbf{condottiero} (qui Micheletto Attendolo), con una ``casa'' e una corte
piuttosto vaste: il cuore militare e materiale della compagnia. La compagnia si
sposta su tutto il territorio italiano e \textbf{cambia committente}: Micheletto
combatté per tutti i principali stati del primo Quattrocento (Regno di Napoli,
Stato della Chiesa, Repubblica di Venezia, Repubblica di Firenze, Stato di
Milano), cambiando parte ``come era normale fare''. Era un aristocratico, dei
conti di Cotignola, e possedeva diversi feudi: un'organizzazione complessa.

Dentro la compagnia c'erano molti altri condottieri, piccoli e grandi. L'unità
militare di base era la \textbf{lancia}, composta da \textbf{tre uomini}: il
cavaliere pesantemente armato, un servente e un ragazzo che teneva i
cavalli\footnote{La trascrizione è confusa (``un servente e il tiro, un
servitore che teneva in giro''); la composizione a tre della lancia è quella
indicata dal relatore.}. Un condottiero poteva avere da una sola lancia a
centinaia di uomini.

% ---------------------------------------------------------------------
\section{Il cortocircuito: dai libri contabili al gioco di ruolo}
\label{sec:c1-2-cortocircuito}

Prima di studiare storia a Siena, Belsanti era già un \textbf{nerd} (``il
termine esisteva già'', ma non era ancora sdoganato): appassionatissimo
giocatore di \textbf{giochi strategici} e di \textbf{giochi di ruolo}, aveva
fatto il master, ``smontava sistemi, creava campagne''. Il contatto con le
fonti produsse un \textbf{cortocircuito immediato}.

I libri contabili sembrano qualcosa di freddo, solo numeri. Ma il tesoriere
\textbf{annotava spesso gli eventi} che accadevano ai condottieri: c'era in lui
``la passione non solo di fare i conti, ma anche di raccontare''. Ne emergono
personaggi ``incredibili'' (Squarcione, un mulattiere dal nome pittoresco, mezzi
frati, rinnegati: ``una variegatissima umanità'') e anche un conterraneo del
relatore, \textbf{Cola da Trani}\footnote{Trascritto ``Col Vitrani'' e ``Con la
Vitrana''. Trani è la città pugliese ``famosa per la sua cattedrale sul
mare''.}.

I condottieri seguono l'intero arco di vita della compagnia, che dura decine di
anni, e dai libri si ricostruisce \textbf{la loro carriera}: molti partono con
pochi uomini e, man mano che la compagnia ottiene ingaggi migliori, ne
acquistano di più, ``esattamente come un classico \emph{upgrade} di un RPG''.
Cola da Trani attraversa così quarant'anni di un'epoca di grandi conflitti,
prima che la \textbf{pace di Lodi} (1454\footnote{Il relatore dice ``1453'';
la pace di Lodi fu firmata nell'aprile 1454.}) consolidi gran parte dei confini
fino all'Unità d'Italia. Poi, un giorno, la fine: la nota del 7 giugno racconta
che Cola da Trani muore \textbf{affogato nell'Oglio}, a Brescia, ``per pigliare
un'anguilla in domenica'', trascinato giù dalla pesante armatura che non si era
tolto. Una storia ``tragicomica'' che lo fece ``viaggiare immediatamente con
l'immaginazione''.

Nei libri, accanto ai nomi, ci sono i \textbf{numeri delle lance} di ciascun
condottiero e talvolta dei \textbf{cavalli} al seguito (da cui Belsanti ha
ricavato le lance al contrario: un condottiero con 13 cavalli ha quattro lance
più un cavallo). La fanteria era organizzata diversamente, e cominciavano a
comparire gli \textbf{schioppettieri}, con le prime armi da fuoco, e i grandi
scudi \textbf{pavesi}\footnote{Trascritto ``scoppiettieri'' e ``scudi
palvezzati''.}.

% ---------------------------------------------------------------------
\section{Il sogno interrotto e la nascita dello studio}
\label{sec:c1-2-studio}

Dopo la tesi, il piano era fondare con l'università una società per la
\textbf{multimedializzazione della storia}: digitalizzazione delle fonti ma anche
divulgazione, quella che poi si sarebbe chiamata \textbf{Public History}. Un buon
gruppo di ragazzi era pronto. Ma accaddero due cose: una ``guerra
interdipartimentale orribile'' da cui il suo professore uscì sconfitto, e poco
dopo un trauma finanziario dell'Università di Siena. Belsanti non volle restare
a fare il ricercatore precario e tornò a casa.

Circa un anno dopo (``sette mesi, in realtà'') incontrò ``un gruppo di pazzi''
che stava sviluppando un videogioco senza alcun legame con la sua formazione e
aveva bisogno di qualcuno che, come un condottiero, li guidasse. Non sapeva
nulla di impresa (``per me una fattura era quella che facevano le streghe'') e
si mise a studiare. È qui che, dice, la formazione umanistica si è rivelata ``la
cosa più utile'': l'università italiana, quella di Siena, i suoi professori gli
hanno dato la capacità di affrontare \textbf{qualsiasi tema}, di costruire
bibliografie e di studiare. Critica i sistemi educativi che non riescono a far
passare la conoscenza come passione e non come competizione per i voti: la
storia non è un dato, sono ``storie di altri esseri umani''.

Lo studio, fin dall'inizio, ha avuto una linea \textbf{storico-culturale ed
educativa}. Cita alcuni progetti educativi:
\begin{itemize}
  \item con la \textbf{Regione Puglia}, sugli \emph{eco-warriors} per il settore
        ambientale\footnote{Trascritto ``eco-worners''.};
  \item con l'\textbf{INAIL}, sulla comunicazione della \textbf{cultura della
        sicurezza} a casa, a scuola e sui luoghi di lavoro: oltre
        \textbf{100.000 studenti} coinvolti e ``scuole entusiaste''.
\end{itemize}
Lo studio ha attraversato ``tantissime difficoltà'' con istituzioni che all'epoca
non sapevano nulla né di videogiochi né di imprenditoria giovanile (``avremmo
dovuto fare a cazzotti''). Il sogno della tesi restò dormiente.

\subsection{Videogames e alta cultura}

Tra il 2017 e il 2018 l'incontro con un professore\footnote{Il nome è
trascritto ``Giovanni Stracci'' e non è verificabile.} riaccese il progetto e
lo spinse a pubblicare la tesi online. Nel 2018 Belsanti fondò un evento
chiamato \textbf{Videogames e Alta Cultura} (\emph{Videogames and High Culture}),
un progetto internazionale con filosofi, storici e accademici, in cui ``alta
cultura'' è intesa in senso nuovo: interrogarsi su che cosa significhi oggi
cultura. L'evento ottenne subito un riconoscimento da \textbf{Gamescom} in
Germania (attraverso l'area oggi chiamata \textbf{devcom}\footnote{Trascritto
``Tentacom'' e ``Diviscon Dev''.}), con cui sono stati realizzati \textbf{otto
spin-off consecutivi}; a breve ne arriverà una nuova edizione.

Le sessioni indagano la storia del gioco in senso lato: il videogioco è solo una
``parentesi'' della grande storia del gioco che accompagna la civiltà, a partire
dalla tesi di \textbf{Huizinga}\footnote{Trascritto ``Osiga''.}: non solo i
giochi fanno parte della cultura, ma \textbf{la cultura stessa nasce nel gioco}.

% ---------------------------------------------------------------------
\section{Il game design come lente sul mondo}
\label{sec:c1-2-lente}

Belsanti insegna game design dal 2020, anche in università a Roma\footnote{Il
nome dell'ateneo è illeggibile (``da Rincarinus'').}, e di recente ha avviato in
Puglia il primo corso pubblico del settore presso un ITS (Apulia
Digital\footnote{``IKS a Puglia Digital'' nella trascrizione: con ogni
probabilità un Istituto Tecnico Superiore pugliese.}). Il suo punto di partenza
è che il game design \textbf{non riguarda solo la professione} di game designer
e la costruzione dei giochi, ma è \textbf{una lente fondamentale per indagare
il passato, il presente e il futuro}.

Ama provocare dicendo che \textbf{l'intero sistema umano della civilizzazione e
delle leggi è un regolamento}, scritto spesso male, caotico, a volte opprimente;
uno dei motivi per cui amiamo i giochi è che i loro regolamenti sono molto più
equi\footnote{La trascrizione dice ``molto più fedeli dalla realtà in cui siamo
costretti a leggere'': il senso sembra ``più chiari ed equi delle regole della
realtà in cui siamo costretti a vivere''.} di quelli della realtà.

% ---------------------------------------------------------------------
\section{Il primo documento di progetto (2022)}
\label{sec:c1-2-documento}

Lo sviluppo vero comincia \textbf{verso la fine del 2022}, con la
partecipazione al bando \textbf{Creative Europe MEDIA}\footnote{Trascritto
``Creative Europe FIGE'' e ``Creative India''.}, ``il più importante fondo
europeo per lo sviluppo dei videogiochi'', vinto con un piccolo contributo.
Belsanti mostra il documento con cui fu vinto il finanziamento:
\begin{description}
  \item[Scheda di progetto.] Titolo, informazioni generali e una
        ``ottimistica'' \emph{target release}, ``che viene sempre cambiata''.
  \item[Target.] Riferendosi a \emph{The Art of Game Design} di Jesse
        Schell\footnote{Trascritto ``gestis shell''.}, uno dei libri base del suo
        corso, osserva che un'azienda deve sempre capire se esiste un pubblico
        sufficiente. È convinto che esistano abbastanza giocatori appassionati
        di un \textbf{gioco strategico} con discussioni anche filosofiche sul
        senso stesso del conflitto; ma raggiungerli è difficilissimo in un mercato
        ``affogato di progetti''. Ammette che il progetto è nato ``da
        indipendente folle'': la priorità non è mai stata vendere milioni di
        copie né il marketing aggressivo, ma il sentire ``che fosse giusto che
        quest'opera nascesse''.
  \item[Tetrade elementare.] Il classico schema delle componenti fondamentali
        di un gioco (estetica, meccaniche, storia, tecnologia). Per
        l'\textbf{estetica}, il punto di riferimento esterno forte è la
        Battaglia di San Romano, con qualche anacronismo voluto: il gioco si
        spinge con alcuni elementi fino al Cinquecento.
\end{description}

\subsection{Un videogioco non insegna la storia}

Qui Belsanti fa una dichiarazione di principio, ``proprio in quanto storico'':
è \textbf{profondamente contrario} agli slogan di marketing secondo cui un
videogioco o un romanzo storico \emph{insegnano} la storia. Citando uno
studioso del romanzo storico\footnote{Il nome è illeggibile (``Antiano
Chiarini'').}: \textbf{la favola vince sulla ricostruzione storica}. Si può
cercare di essere il più accurati possibile, ``ma non sto insegnando storia. La
storia è un'altra cosa''. Un'opera ludica o un romanzo possono affascinare,
portare conoscenze e spingere ad approfondire, ma le due cose \textbf{vanno
distinte}.

Detto questo, il gioco ha un primato di cui va fiero: per quanto gli risulta, è
\textbf{il primo videogioco al mondo che usa fonti storiche archivistiche
inedite come base narrativa}. È una scommessa su un certo tipo di giocatori
``nuovi, adulti, complessi, vecchi come me''.

\subsection{Le componenti del progetto}

\begin{description}
  \item[Meccaniche.] Uno \textbf{strategico} con fasi di \emph{wargame} molto
        intense ``per chi vorrà giocarle'', e in alternativa battaglie
        \textbf{narrative} semplificate.
  \item[Storia.] Un elemento narrativo molto forte, impostato come una
        \textbf{tragedia}. Una scelta particolare riguarda la lingua: oltre a
        italiano e inglese, si potrà giocare in \textbf{italiano antico del
        Quattrocento}, costruito da un teorico della letteratura amico dello
        studio (un lavoro complesso, con molti dialetti; ``scherzando'' lo hanno
        chiamato ``rinascimentale''), e in un \textbf{inglese arcaico}
        ``shakespeariano''.
  \item[Loop di gioco.] Tutti i giochi sono fatti di \textbf{azioni
        ricorsive}. Il loop principale alterna tre attività: le
        \textbf{battaglie}, lo \textbf{svolgimento della storia} e la parte
        \textbf{manageriale}, più una serie di azioni secondarie.
\end{description}

Seguono i primi studi grafici: personaggi ripresi da ritratti di condottieri
dell'epoca; un \emph{mock-up} dell'\textbf{accampamento} (la fase narrativa) con
uno sfondo ricavato da un quadro di Paolo Uccello; le \textbf{unità militari}
sulla mappa, in parte prese direttamente dal dipinto e in parte ricostruite da
artisti con tecnica simile.

\subsection{Un wargame esagonale}

Le prime immagini delle meccaniche mostrano una \textbf{mappa esagonale} e un
gioco \textbf{a turni}. Belsanti si definisce appassionato di una ``nicchia della
nicchia'', il \emph{wargame} da tavolo: ha amato il sistema \emph{World in
Flames}\footnote{Trascritto ``Word in Frames''. Il relatore aggiunge di esserne
stato ``campione'' nel 1996 (``campione svedese'' nella trascrizione,
probabilmente ``campione italiano'') e lo descrive come un gioco che dura in
media ``150'' ore con ``14.000 pedine di cartone''. Le cifre sono riportate come
trascritte.} e una parte di quella ``follia'' è finita nella simulazione:
\begin{itemize}
  \item ogni unità sull'esagono ha \textbf{sei fronti} ed è orientata: fronte,
        fianchi e retro (attaccare alle spalle è più efficace che frontalmente);
  \item il motore di battaglia è \textbf{proprietario} e verrà riusato in un
        progetto futuro;
  \item ogni unità ha una lunga lista di caratteristiche: lance, età, attacco a
        distanza, attacco in mischia e \textbf{lealtà}. Quest'ultima viene
        direttamente dalle fonti, dove capita di leggere di condottieri
        \textbf{fuggiti} dall'accampamento dopo essere stati pagati;
  \item c'è un elemento \textbf{epistolare}: si scrivono lettere e si risponde
        alle lettere.
\end{itemize}

% ---------------------------------------------------------------------
\section{Il gioco oggi}
\label{sec:c1-2-gioco}

Il gioco \textbf{uscirà su Steam nel 2027}; la data non è ancora annunciata.
Dopo il prototipo, lo studio ha avuto accesso al \textbf{tax credit} italiano
per lo sviluppo e ha intensificato il lavoro nell'ultimo anno. Nella graduatoria
dei progetti ammessi quell'anno lo studio si è classificato \textbf{al primo
posto}, con un punteggio del 96\%: il tax credit prevede punteggi per il
\textbf{riconoscimento culturale} e bonus per la valorizzazione del patrimonio
artistico italiano, e il primo Rinascimento è ``una di quelle epoche di cui
ancora viviamo di rendita''.

\subsection{Il titolo e il tema}

Il titolo è \emph{Huomo d'arme: The End of Knights}\footnote{Trascritto ``Homo
Titanic the End of Nights'': la ricostruzione, coerente con il personaggio
dell'``uomo d'arme'' descritto subito dopo e con il tema della fine della
cavalleria, è un'ipotesi. Si consiglia di verificare il titolo esatto sulla
pagina Steam del gioco.}. Il Quattrocento, e i libri contabili lo mostrano bene,
è un'epoca di grande trasformazione, l'\textbf{autunno del Medioevo}: la
\textbf{cavalleria} è ormai soprattutto un titolo onorifico, slegato dall'origine
militare e dal merito sul campo; gli \textbf{uomini d'arme}, invece, sono ciò
che resta di quel merito. Non sono necessariamente nobili, anzi spesso non lo
sono; ma nelle fonti compare la formula con cui qualcuno ``fu fatto uomo d'arme
in nome di Dio e di San Giorgio'', con una ritualità ripresa dal
\textbf{sogno cavalleresco}: il desiderio, sospeso tra realtà e immaginazione,
del guerriero custode della giustizia. Intanto si affaccia sempre di più il
mondo \textbf{borghese e mercantile}. Su questo passaggio d'epoca si concentra
la narrazione.

\subsection{L'interfaccia e la narrativa}

\begin{itemize}
  \item Il \textbf{menu} è ``molto esoterico'': un archivio con icone sospese
        nell'aria. Una conduce alle \textbf{fonti storiche}, che si potranno
        leggere (la tesi è in traduzione in inglese, accelerata dagli strumenti
        di IA).
  \item Il gioco racconta una \textbf{tragedia shakespeariana} in cui i
        personaggi attraversano \textbf{quattro epoche} della loro vita, legate
        alle \textbf{quattro stagioni}.
  \item Al centro c'è il protagonista, ispirato a una delle figure più
        misteriose dei libri contabili: una sorta di ``legionario straniero''
        che \textbf{non viene mai chiamato per nome}, ma solo ``uomo d'arme''.
        È sempre in armatura e non mostra mai il volto.
  \item Accanto a lui i compagni: Iacopo, Scorciola e altri; una figura
        femminile forte, \textbf{Polissena}, figlia di un condottiero e
        rappresentata anch'essa in armatura (nelle fonti compare una Polissena,
        e Belsanti ha ``navigato'' anche con questa storia); un
        \textbf{giullare} che scherza ``in modo molto serio'' sul concetto
        stesso di guerra.
  \item Nelle fasi narrative ci sono \textbf{scelte etiche}.
\end{itemize}

\subsection{I quattro indicatori}

La compagnia è descritta da \textbf{quattro grandi indicatori}:
\begin{enumerate}
  \item il \textbf{tesoro}, in fiorini (``il dollaro del Medioevo'');
  \item il \textbf{prestigio aristocratico};
  \item la \textbf{gloria militare};
  \item il \textbf{sogno cavalleresco}.
\end{enumerate}
Il cuore del design è la \textbf{tensione} fra di essi: se si cerca di preservare
il sogno cavalleresco, spesso si entra in conflitto con tutti gli altri.

\subsection{Mappa, gestione ed esercito}

\begin{itemize}
  \item Una mappa d'Italia \textbf{volutamente imprecisa} dal punto di vista
        geografico, che ricostruisce l'assetto del Quattro--Cinquecento: gli
        stati offrono \textbf{condotte} (contratti) e il giocatore sceglie dove
        andare a combattere, cambiando committente nel corso della partita.
  \item Una fase \textbf{gestionale}: il condottiero gestisce i propri
        \textbf{feudi}, che può comprare o vendere. ``Non bisogna mai andare in
        rosso'': il condottiero è \textbf{un imprenditore} e deve stare molto
        attento ai costi.
  \item L'\textbf{esercito}: ogni unità ha molte caratteristiche e può ricevere
        \textbf{upgrade}. Ma le unità che oggi fanno parte del tuo esercito
        domani possono essere \textbf{nell'esercito avversario}: potenziarle
        significa rischiare di ritrovarsele contro. Per lo stesso motivo, nel
        mondo del mercenariato, non ci si uccide ``malamente'': si fanno
        \textbf{prigionieri e riscatti}, perché ``da un momento all'altro ci
        potrebbe essere un cambio di casacca''.
\end{itemize}

\subsection{Due modi di combattere}

\begin{description}
  \item[Battaglia narrativa.] Per chi non vuole affrontare il \emph{wargame},
        una battaglia molto semplificata su \textbf{tre fronti}, preceduta da un
        momento di dialogo.
  \item[Battaglia tattica.] Il sistema \textbf{esagonale}, da \emph{board game}
        con miniature e schieramenti, con tutte le regole citate e la
        \textbf{zona di controllo} (\emph{zone of control}) tipica dei wargame:
        l'area intorno a un'unità che limita i movimenti liberi delle unità
        nemiche. ``Chi ama questi giochi capisce'' e apprezzerà la profondità;
        Belsanti spera anche che qualcuno la scopra e impari ad
        apprezzarla.
\end{description}

Infine, il gioco conterrà una \textbf{raccolta delle fonti storiche} con accesso
completo alla tesi. Il trailer non si riesce a mostrare con l'audio, e il relatore
apre alle domande.

% ---------------------------------------------------------------------
\section{Le domande del pubblico}
\label{sec:c1-2-domande}

\subsection{Il compromesso tra storico e autore}

La prima domanda viene dalla \textbf{curatrice del museo}, anch'essa
storica\footnote{Si presenta come ``Caterina Pasca''; il cognome non è
verificabile.}: che compromessi ha dovuto fare, da storico, nel trasporre ciò che
ha studiato in un'opera artistica, che deve cedere il passo alla drammaturgia?
Ci sono cose sacrificate all'immaginario comune sul tardo Medioevo?

Belsanti dice di aver affrontato la questione ``in modo pluriennale'' e di essere
arrivato alla consapevolezza che \textbf{l'accuratezza storica, in un'opera
narrativa, è tra virgolette una bufala}. Non voleva però le ``storpiature'' di
certe opere (cita un film su Ercole dei primi anni Duemila e \emph{Dante's
Inferno}). Ha fatto ``molta pace con se stesso'': dove ha inserito elementi
storici accurati, lo sono \textbf{pienamente}; i personaggi della compagnia
\textbf{sono nomi che stanno nelle fonti}, e chi vuole può verificarli
direttamente sulla tesi e sulle scansioni (``saranno pochissimi, come saranno
pochissimi quelli che giocheranno in italiano rinascimentale, ma era giusto
così''). Nella scheda del progetto lo studio ha dichiarato apertamente le scelte
\textbf{anacronistiche} (il fulcro è il primo Quattrocento, ma ci si spinge fino
al Cinquecento, come hanno fatto anche i linguisti): \textbf{dichiararle} è per
loro molto importante.

Richiama poi il professor Antonio Brusa\footnote{Trascritto ``Antonio Cusa''.
Antonio Brusa è uno storico della didattica, tra i promotori dell'associazione
Historia Ludens (trascritta ``sono fondati su Le Monde''); l'identificazione è
un'ipotesi.}, pioniere della didattica della storia con i giochi: quando gli
insegnanti usano i giochi, la cosa più importante non è l'accuratezza del
singolo dettaglio, ma ricondurre a \textbf{dinamiche plausibili} e far capire i
\textbf{meccanismi della storia}. Se il giocatore capisce che una compagnia di
ventura è \textbf{un'attività imprenditoriale} in cui è vitale non andare in
rosso, ``hai già vinto'', anche se la rendita attribuita a un piccolo feudo
della campagna fiorentina fosse sbagliata.

\subsection{Il debriefing}

Aggiunge un punto che considera centrale: i giochi che vogliono essere educativi
hanno bisogno di qualcosa che i videogiochi d'intrattenimento non potranno
``praticamente mai'' avere, il \textbf{debriefing}. Dopo aver giocato bisogna
parlare con lo storico, con il docente. L'esperienza narrativa è coinvolgente
ed emotiva, e \textbf{l'emozione è un po' nemica della scienza}: bisogna entrarci
ma poi anche \textbf{uscirne}. I videogiochi sono fatti per catturare
l'attenzione all'infinito (``l'economia dell'attenzione''); per fare davvero
ricerca storica serve un'uscita. Il suo ``piccolo debriefing'' è dare accesso
alle fonti: ``se vuoi, là c'è tutto lo storico; questa invece è una
creazione''.

Conclude che non ha fatto grandi compromessi: non gli è stato difficile mettere
da parte la voce dello storico per dar voce allo \textbf{scrittore e
sceneggiatore}. Anche gli storici, del resto, finiscono spesso per
``innamorarsi'' delle persone che studiano, per un meccanismo di
identificazione narrativa; lui si è lasciato andare, sapendo di poter poi
riprendere le distanze.

\subsection{Trent'anni per arrivare alla visione finale}

Un secondo spettatore chiede se, dalla tesi del '98 al primo bando, ci sia stata
la tentazione di pubblicare subito il gioco com'era nel primo prototipo, invece
di arrivare dopo quasi trent'anni a una visione finale così curata.

La risposta è che l'opera si è concretizzata ``abbastanza rapidamente, a un
certo punto'', anche per via delle scadenze. Trent'anni fa c'era solo il
desiderio di dar vita a un'esperienza da condottiero; ``fosse stato per me
sarebbe rimasta un'eterna incompiuta''. Ammette la paura del confronto con un
pubblico moderno ``affogato di opere'', spesso superficiale e cattivo, e una sorta
di ``schizofrenia artistica'': voler dire delle cose e insieme no; e un
``bisogno cavalleresco'' di lasciare l'opera come tensione incompiuta. Sono stati
i finanziamenti a forzare la concretizzazione: il bando Creative Europe chiedeva
solo un prototipo, ma il \textbf{tax credit obbliga a uscire} (almeno all'epoca;
ora le regole sono cambiate).

L'idea iniziale era più ambiziosa: un \textbf{simulatore di compagnia di
ventura} libero e aperto, con un condottiero qualsiasi e personalizzabile. Poi,
artisticamente, ha scelto una connotazione \textbf{narrativa più piccola e più
forte}: un'esperienza di \textbf{4--5 ore}, molto concentrata, con la struttura
della \textbf{tragedia shakespeariana} (Shakespeare è uno dei suoi autori
preferiti). ``Il fatto di essere stato costretto ad accelerare ha prodotto
qualcosa di diverso, spero unico.'' A prescindere dal risultato commerciale e
critico, sarà contento di averlo \textbf{terminato}.

\subsection{La colonna sonora}

L'ultima domanda è di una spettatrice venuta da Milano\footnote{Si presenta
come ``Yolanda, di Minuto''; il nome non è verificabile.}, interessata alle
colonne sonore dei videogiochi, che spesso vivono di vita propria: ci sarà
musica, e quale, in un gioco in cui il periodo storico è fondamentale?

Belsanti racconta una scelta ``molto forte, totalmente antagonistica''. Il
suono ambientale è curato da un sound designer, e nella fase narrativa
dell'accampamento ci sono gli effetti delle stagioni (primavera, inverno,
pioggia\dots). Ma poiché il gioco è diviso in \textbf{quattro stagioni}, per le
battaglie ha deciso di usare \textbf{Le quattro stagioni di Vivaldi}, in una
versione dell'etichetta Brilliant Classics\footnote{Trascritto ``Brilliant
Plastics''.} non centrata sul violino solista ma su un \textbf{flauto}. Nel
trailer e nell'epilogo ci sono anche passi del \textbf{Requiem di
Mozart}\footnote{Trascritto ``il record di Mozart''.}. È una scelta di forte
\textbf{anacronismo}, di separazione dalla ricostruzione storica delle musiche
del Quattro--Cinquecento (più difficile da recuperare per il primo Quattrocento,
più documentata verso il Cinquecento), motivata dall'asse narrativo delle
quattro stagioni e dal contrasto che crea. Il tentativo di mostrare il trailer
con l'audio fallisce di nuovo.

% ---------------------------------------------------------------------
\section{La premiazione}
\label{sec:c1-2-premio}

L'intervento si chiude con un premio a Belsanti ``per il suo ruolo di
\textbf{pioniere} nello sviluppo dell'industria videoludica italiana, per aver
saputo coniugare storia, cultura e game design, contribuendo a dimostrare il
valore culturale, artistico ed espressivo delle opere interattive''. Gli
organizzatori si scusano per gli inconvenienti tecnici e si accordano per
spedirgli il premio.

% ---------------------------------------------------------------------
\section*{Punti chiave del capitolo}
\addcontentsline{toc}{section}{Punti chiave del capitolo}

\begin{itemize}
  \item Un videogioco strategico nato da una \textbf{tesi di laurea del 1998}
        sui libri contabili della compagnia di ventura di Micheletto Attendolo
        (archivio della Fraternita dei Laici di Arezzo), con la
        \textbf{Battaglia di San Romano} di Paolo Uccello come riferimento
        estetico.
  \item Le fonti contabili contengono \textbf{storie}: carriere di condottieri
        che crescono ``come un upgrade di un RPG'', morti tragicomiche (Cola da
        Trani affogato nell'Oglio per un'anguilla), fughe dopo la paga.
  \item Il game design è \textbf{una lente sul mondo}: la civiltà stessa è un
        regolamento, spesso scritto male; la cultura nasce nel gioco
        (Huizinga).
  \item Finanziamenti: \textbf{Creative Europe MEDIA} per il prototipo (2022),
        poi il \textbf{tax credit} italiano (primo in graduatoria, 96\%), che
        obbliga a pubblicare. Uscita su Steam prevista nel 2027.
  \item Un videogioco \textbf{non insegna la storia}: ``la favola vince sulla
        ricostruzione storica''. Le scelte anacronistiche vanno
        \textbf{dichiarate}; ciò che conta è la \textbf{plausibilità delle
        dinamiche}. Per un uso educativo serve il \textbf{debriefing}.
  \item Design: loop su battaglie, storia e gestione; quattro indicatori
        (tesoro, prestigio aristocratico, gloria militare, \textbf{sogno
        cavalleresco}) in tensione tra loro; unità che possono passare al
        nemico; lealtà ricavata dalle fonti; wargame esagonale con fronti e
        zone di controllo oppure battaglia narrativa semplificata.
  \item Ridurre l'ambizione (dal simulatore aperto a una tragedia di 4--5 ore)
        ha dato \textbf{un'identità più forte}. Musica: scelta volutamente
        anacronistica (Vivaldi, Mozart) legata alle quattro stagioni.
\end{itemize}
